% 6thGradeReviewExplanations — pages 34–42
% Section K (K1–K11) and Section L (L1).

\lesson{K. Expressions, Factors, And Fraction Structure}

\question{K1}

\blockhead{Question}
Is $8x + 1$ a linear expression?

\checked{Yes}

\blockhead{What the question is asking}
This question asks whether $8x + 1$ has the letter $x$ only to the first power (a small raised 1; no $x^{2}$, and no letter under a top-over-bottom bar).

\blockhead{Important words}
\begin{words}
  \item An \term{expression}is a math phrase with numbers, variables, and operations (no equals sign required).
  \item A \term{linear expression}has variables only to the first power. There is no $x^{2}$, no $x^{3}$, and no variable in a denominator. Typical form: $ax + b$.
  \item \term{First power}means a small raised 1 on the variable (written $x^{1}$, or just $x$).
  \item A \term{constant}is a number that stands alone with no variable letter (like 1 or 8).
  \item The term $8x$ means $8 \times x$, and $x$ is to the first power.
\end{words}

\blockhead{Work the problem one tiny step at a time}
\begin{steps}
  \item Look at each term in $8x + 1$.
  \item The term $8x$ has variable $x$ to the power 1.
  \item The term 1 is a constant (no variable).
  \item There is no $x^{2}$ or higher power.
  \item So $8x + 1$ is a linear expression. Answer: Yes.
\end{steps}

\finalans{Yes}

\blockhead{Why this makes sense}
Contrast with $8x^{2} + 1$, which is not linear because of $x^{2}$. Here the power of $x$ is 1, so it is linear.

\question{K2}

\blockhead{Question}
Prime factorization of $90 = $

\checked{$2 \cdot 3^{2} \cdot 5$}

\blockhead{What the question is asking}
This question asks you to write 90 as a times-chain of numbers that each can be divided evenly only by 1 and themselves.

\blockhead{Important words}
\begin{words}
  \item A \term{factor}of a number divides that number evenly (no leftover).
  \item A \term{prime number}has exactly two factors: 1 and itself. Examples: $2, 3, 5, 7, 11$.
  \item \term{Prime factorization}writes a number as a product of primes only.
  \item An \term{exponent}like $3^{2}$ means $3 \times 3$.
\end{words}

\blockhead{Work the problem one tiny step at a time}
\begin{steps}
  \item Factor 90 into any pair: $90 = 9 \times 10$, or better, start dividing by small primes.
  \item 90 is even, so divide by 2: $90 = 2 \times 45$.
  \item $45 = 5 \times 9$, and $9 = 3 \times 3$.
  \item So $90 = 2 \times 3 \times 3 \times 5$.
  \item Write repeated factors with an exponent: $2 \cdot 3^{2} \cdot 5$.
\end{steps}

\finalans{$2 \cdot 3^{2} \cdot 5$}

\blockhead{Why this makes sense}
Check by multiplying: $3^{2} = 9$, then $2 \times 9 = 18$, then $18 \times 5 = 90$. Correct.

\question{K3}

\blockhead{Question}
Find the GCF of $27$ and $36$:

\checked{9}

\blockhead{What the question is asking}
This question asks for the biggest whole number that divides both 27 and 36 evenly (no leftover).

\blockhead{Important words}
\begin{words}
  \item A \term{factor}of a number divides that number evenly.
  \item A \term{common factor}divides both numbers.
  \item The \term{GCF}(greatest common factor) is the largest common factor.
  \item To \term{prime factorize}means write the number as a product of prime numbers only.
  \item \term{Shared primes}are prime factors that appear in both prime factorizations.
\end{words}

\blockhead{Work the problem one tiny step at a time}
\begin{steps}
  \item Prime factorize 27: $27 = 3 \times 3 \times 3 = 3^{3}$.
  \item Prime factorize 36: $36 = 2 \times 2 \times 3 \times 3 = 2^{2} \cdot 3^{2}$.
  \item Shared primes: only 3 appears in both; the smaller power is $3^{2}$.
  \item So GCF $= 3^{2} = 9$.
  \item List check: factors of 27: $1, 3, 9, 27$. Factors of 36: $1, 2, 3, 4, 6, 9, 12, 18, 36$. Largest shared: 9.
\end{steps}

\finalans{9}

\blockhead{Why this makes sense}
9 divides 27 ($27 \div 9 = 3$) and divides 36 ($36 \div 9 = 4$). Nothing larger than 9 does both.

\question{K4}

\blockhead{Question}
Simplify $\dfrac{27}{36} = $

\checked{$\dfrac{3}{4}$}

\blockhead{What the question is asking}
This question asks you to shrink $\tfrac{27}{36}$ by dividing the top and bottom by the same whole number until nothing bigger than 1 divides both evenly.

\blockhead{Important words}
\begin{words}
  \item To \term{simplify}a fraction, divide top and bottom by their GCF.
  \item \term{Simplest terms}means the top and bottom share no common factor greater than 1.
  \item From K3, the GCF of 27 and 36 is 9.
\end{words}

\blockhead{Work the problem one tiny step at a time}
\begin{steps}
  \item From K3, the GCF of 27 and 36 is 9 (shared primes give $3^{2} = 9$).
  \item Divide the numerator by 9: $27 \div 9 = 3$.
  \item Divide the denominator by 9: $36 \div 9 = 4$.
  \item Write $\dfrac{3}{4}$.
  \item Check: 3 and 4 share no common factor greater than 1, so it is simplest.
\end{steps}

\finalans{$\dfrac{3}{4}$}

\blockhead{Why this makes sense}
$\dfrac{27}{36}$ and $\dfrac{3}{4}$ name the same amount. Cross check: $27 \times 4 = 108$ and $36 \times 3 = 108$.

\question{K5}

\blockhead{Question}
Convert $\mixed{2}{2}{3}$ to an improper fraction:

\checked{$\dfrac{8}{3}$}

\blockhead{What the question is asking}
This question asks you to rewrite $\mixed{2}{2}{3}$ (a whole number plus equal parts over a whole) as one top-over-bottom amount where the top is at least as big as the bottom.

\blockhead{Important words}
\begin{words}
  \item A \term{mixed number}has a whole number and a fraction, like $\mixed{2}{2}{3}$.
  \item An \term{improper fraction}has a numerator greater than or equal to the denominator, like $\dfrac{8}{3}$.
  \item Rule: improper numerator $= $ (whole $\times$ denominator) $+$ numerator. Keep the same denominator.
\end{words}

\blockhead{Work the problem one tiny step at a time}
\begin{steps}
  \item Whole number $= 2$, numerator $= 2$, denominator $= 3$.
  \item Multiply whole by denominator: $2 \times 3 = 6$.
  \item Add the old numerator: $6 + 2 = 8$.
  \item Keep denominator 3: $\dfrac{8}{3}$.
\end{steps}

\finalans{$\dfrac{8}{3}$}

\blockhead{Why this makes sense}
Check: $\tfrac{8}{3} = 2 + \tfrac{2}{3}$ because $\tfrac{6}{3} = 2$ and $\tfrac{2}{3}$ remains. Same mixed number.

\question{K6}

\blockhead{Question}
Which is greater: $\mixed{1}{3}{4}$ or $(1)\!\left(\tfrac{3}{4}\right)$?

\checked{$\mixed{1}{3}{4}$}

\blockhead{What the question is asking}
This question asks which value is larger: $\mixed{1}{3}{4}$ (one and three fourths) or $(1)\!\left(\tfrac{3}{4}\right)$ (one times three fourths).

\blockhead{Important words}
\begin{words}
  \item $\mixed{1}{3}{4}$ means $1 + \dfrac{3}{4}$.
  \item $(1)\!\left(\tfrac{3}{4}\right)$ means 1 times $\dfrac{3}{4}$, which is just $\dfrac{3}{4}$.
  \item Parentheses next to each other mean multiply.
\end{words}

\blockhead{Work the problem one tiny step at a time}
\begin{steps}
  \item Rewrite the mixed number: $\mixed{1}{3}{4} = 1 + \dfrac{3}{4} = \dfrac{4}{4} + \dfrac{3}{4} = \dfrac{7}{4}$.
  \item Compute the product: $(1)\!\left(\tfrac{3}{4}\right) = \dfrac{3}{4}$.
  \item Compare $\dfrac{7}{4}$ and $\dfrac{3}{4}$. Same denominator, so compare tops: $7 > 3$.
  \item Therefore $\mixed{1}{3}{4}$ is greater.
\end{steps}

\finalans{$\mixed{1}{3}{4}$}

\blockhead{Why this makes sense}
A mixed number includes a whole 1 plus a fraction. Multiplying by 1 does not add a whole; it only keeps $\tfrac{3}{4}$. So the mixed number is larger.

\question{K7}

\blockhead{Question}
In $\dfrac{m}{2}$, coeff of $m$ is \blank[1.8cm]

\checked{$\dfrac{1}{2}$}

\blockhead{What the question is asking}
This question asks for the number that multiplies the letter $m$ in $\tfrac{m}{2}$.

\blockhead{Important words}
\begin{words}
  \item A \term{coefficient}is the number multiplying a variable.
  \item $\tfrac{m}{2}$ means $m$ divided by 2, which is the same as $\tfrac{1}{2} \times m$.
\end{words}

\blockhead{Work the problem one tiny step at a time}
\begin{steps}
  \item Rewrite $\dfrac{m}{2}$ as $\dfrac{1}{2}m$.
  \item The number multiplying $m$ is $\dfrac{1}{2}$.
  \item So the coefficient of $m$ is $\dfrac{1}{2}$.
\end{steps}

\finalans{$\dfrac{1}{2}$}

\blockhead{Why this makes sense}
Dividing by 2 is the same as multiplying by one-half. The coefficient is not 2; it is $\tfrac{1}{2}$.

\question{K8}

\blockhead{Question}
In $\dfrac{q}{5}$, coeff of $q$ is \blank[1.8cm]

\checked{$\dfrac{1}{5}$}

\blockhead{What the question is asking}
This question asks for the number that multiplies the letter $q$ in $\tfrac{q}{5}$.

\blockhead{Important words}
\begin{words}
  \item A \term{variable}is a letter that stands for a number. Here the variable is $q$.
  \item A \term{coefficient}is the number that multiplies a variable.
  \item A \term{fraction}$\tfrac{q}{5}$ means $q$ divided by 5.
  \item Dividing by 5 is the same as multiplying by $\tfrac{1}{5}$: $\tfrac{q}{5} = \tfrac{1}{5} \times q$.
\end{words}

\blockhead{Work the problem one tiny step at a time}
\begin{steps}
  \item Start with $\dfrac{q}{5}$.
  \item Rewrite dividing by 5 as multiplying by $\dfrac{1}{5}$: $\dfrac{q}{5} = \dfrac{1}{5}q$.
  \item The number multiplying $q$ is $\dfrac{1}{5}$.
  \item So the coefficient is $\dfrac{1}{5}$.
  \item Numeric check: if $q = 10$, then $\dfrac{10}{5} = 2$ and $\dfrac{1}{5} \times 10 = 2$. Same value, so $\dfrac{1}{5}$ really is the multiplier.
\end{steps}

\finalans{$\dfrac{1}{5}$}

\blockhead{Why this makes sense}
Dividing by 5 means multiplying by $\tfrac{1}{5}$. The coefficient is that multiplier. Check with $q = 10$: both forms give 2.

\question{K9}

\blockhead{Question}
Commutative: $8 + 5 = $

\checked{$5 + 8$}

\blockhead{What the question is asking}
This question asks you to rewrite $8 + 5$ by swapping the order of the two numbers being added.

\blockhead{Important words}
\begin{words}
  \item The \term{commutative property of addition}says $a + b = b + a$. Order can switch; the sum stays the same.
  \item An \term{addend}is a number being added.
\end{words}

\blockhead{Work the problem one tiny step at a time}
\begin{steps}
  \item Start with $8 + 5$.
  \item Switch the order of the two addends.
  \item Write $5 + 8$.
  \item Check: $8 + 5 = 13$ and $5 + 8 = 13$. Same sum.
\end{steps}

\finalans{$5 + 8$}

\blockhead{Why this makes sense}
Commutative means ``order can move.'' Do not change the operation to multiplication; keep addition.

\question{K10}

\blockhead{Question}
Commutative: $7 \times a = $

\checked{$a \times 7$}

\blockhead{What the question is asking}
This question asks you to rewrite $7 \times a$ by swapping the order of the two factors being multiplied.

\blockhead{Important words}
\begin{words}
  \item The \term{commutative property of multiplication}says $a \times b = b \times a$. You may switch the order of the factors; the product stays the same.
  \item A \term{factor}is a number (or variable) being multiplied.
  \item The \term{associative property}only moves parentheses (regroups). It does \textbf{not} switch order. Commutative switches order.
\end{words}

\blockhead{Work the problem one tiny step at a time}
\begin{steps}
  \item Start with $7 \times a$.
  \item Name the factors in order: first factor 7, second factor $a$.
  \item Switch the order of the factors.
  \item After the switch, first factor is $a$, second factor is 7: write $a \times 7$.
  \item Numeric check with a number for $a$: if $a = 3$, then $7 \times 3 = 21$ and $3 \times 7 = 21$. Same product after the order switch.
\end{steps}

\finalans{$a \times 7$}

\blockhead{Why this makes sense}
The product is the same either way. This is not the associative property (associative only regroups parentheses; commutative switches order).

\question{K11}

\blockhead{Question}
Associative: $(b + 3) + 9 = $

\checked{$b + (3 + 9)$}

\blockhead{What the question is asking}
This question asks you to rewrite $(b + 3) + 9$ by moving the curved marks to group the addition in a different way.

\blockhead{Important words}
\begin{words}
  \item The \term{associative property of addition}says $(a + b) + c = a + (b + c)$. The grouping can change; the sum stays the same.
  \item \term{Regroup}means move the parentheses without changing the order of the terms.
  \item An \term{addend}is a number (or letter) being added.
  \item The \term{commutative property}switches order. Associative only moves parentheses; it does not reorder terms.
\end{words}

\blockhead{Work the problem one tiny step at a time}
\begin{steps}
  \item Start with $(b + 3) + 9$. Here $b$ and 3 are grouped first.
  \item List the three addends in order: $b$, then 3, then 9. Keep that order.
  \item Move the parentheses to group 3 and 9: $b + (3 + 9)$.
  \item Numeric check with a number for $b$: if $b = 2$, then $(2 + 3) + 9 = 5 + 9 = 14$.
  \item Same value after regrouping: $2 + (3 + 9) = 2 + 12 = 14$. Check works.
\end{steps}

\finalans{$b + (3 + 9)$}

\blockhead{Why this makes sense}
Associative changes grouping, not order. Writing $9 + (b+3)$ would also equal the same value, but the standard associative rewrite keeps order and only moves parentheses. The numeric check $14 = 14$ confirms the regroup.

\lesson{L. Distributive Property And Factoring}

\question{L1}

\blockhead{Question}
Factor: $8x + 16 = $

\checked{$8(x + 2)$}

\blockhead{What the question is asking}
This question asks you to write $8x + 16$ as a product by pulling out a shared whole number that divides both parts evenly.

\blockhead{Important words}
\begin{words}
  \item To \term{factor}means write a sum as a product (pull out a shared factor).
  \item The \term{GCF}of the terms is the largest factor shared by both terms.
  \item \term{Factored form}means written as a product, like $8(x + 2)$.
  \item Factoring undoes the distributive property: $a(b + c) = ab + ac$.
\end{words}

\blockhead{Work the problem one tiny step at a time}
\begin{steps}
  \item Look at the terms: $8x$ (coefficient 8) and the constant 16.
  \item GCF of 8 and 16 is 8. (Variable $x$ is only in the first term, so it is not part of the GCF.)
  \item Divide first term by 8: $8x \div 8 = x$.
  \item Divide second term by 8: $16 \div 8 = 2$.
  \item Write the factored form: $8(x + 2)$.
  \item Check by distributing: $8 \times x + 8 \times 2 = 8x + 16$. Matches.
\end{steps}

\finalans{$8(x + 2)$}

\blockhead{Why this makes sense}
Always check by distributing back. If you get the original expression, the factoring is correct.
